% !TEX root = ../main.tex
\chapter{Touching the Elephant: Three Views}\label{ch:threeviews}
\begin{paracol}{2}

\begin{sources}
\begin{tabularx}{\linewidth}{@{}l l L@{}}
\toprule
\textbf{Segment} & \textbf{Length} & \textbf{Content}\\
\midrule
\seg{22}{1}{L22.1 FT: Future of banking} & 4:46 & Banking union, cross-border fragmentation, resolution of global banks.\\
\seg{22}{2}{L22.2 Three world views} & 15:14 & Money, economics and finance views; Black's 1976 manifesto.\\
\seg{22}{3}{L22.3 Economics View: Commodity Exchange} & 8:54 & Relative prices, the quantity theory pasted on; Fisher and IS--LM.\\
\seg{22}{4}{L22.4 Finance View: Risk} & 15:41 & CAPM in balance sheets; Black's banking in a world without money.\\
\seg{22}{5}{L22.5 The education of Fischer Black} & 4:52 & ``Noise'' (1985); Treynor's dealer model.\\
\seg{22}{6}{L22.6 Steps from finance view to money view} & 7:45 & Market liquidity, funding liquidity, the central bank.\\
\seg{22}{7}{L22.7 A money view of economics and finance} & 5:12 & Selective control over liquidity allocation.\\
\bottomrule
\end{tabularx}

\smallskip
\textbf{Files.} Transcripts \file{materials/transcripts/L22-P*.txt}; Mehrling's notes \mn{22}{1}.
\end{sources}

\section*{Motivation}
Lecture~1 promised a course that would not be applied microeconomics (utility functions for banks) nor applied macroeconomics (what lies behind the LM curve), but would
start from the world, the actual operations of money markets, and learn inductively \ts{22}{2}{0:03}. The last lecture connects that money view to the economics and
finance that students know: three blind men touching the same elephant.

% ------------------------------------------------------------------
\section{FT: the future of banking}\label{sec:tv-ft}

\begin{casestudy}[FT, December 2012: three data points in one day]
\begin{enumerate}
  \item ``Debate rages on how to set up bank supervisor'': the ECB, with no supervisory experience, no dedicated staff and hostility from some of its future charges, is
        to supervise the eurozone's 6000 banks in a banking union \ts{22}{1}{0:07}. Mehrling's reading of the unspoken resistance: common
        collateral rules would end the sweetheart deals by which national central banks lean on local banks to buy their government's bonds \ts{22}{1}{1:09};
  \item ``Cross-border lending falls sharply'': US banks stop lending to Europe and separate their European and US books; core--periphery lending within Europe shrinks: a
        market-driven ``return to autarky'' \ts{22}{1}{2:10};
  \item an op-ed by Martin Gruenberg (FDIC chairman) and Paul Tucker (Bank of England deputy governor), ``Global banks need global solutions when they fail'': the
        resolution authority would take control of the parent, impose losses on shareholders and unsecured creditors, and remove management. They say they will no longer
        support institutions, not yet what they will do instead \ts{22}{1}{2:51}.
\end{enumerate}
Details from the two articles in the course files (\emph{FT}, 10 December 2012; \file{materials/readings/L22-*.pdf}):
\begin{itemize}
  \item \textbf{Cross-border lending} (Claire Jones, BIS quarterly data). Cross-border claims fell by \$575bn (1.9\%) between the end of March and the end of June 2012, and interbank claims by \$581bn (3.1\%). This was the second sharpest retrenchment since early 2009, ``mainly the result of reduced inter-office positions''. Dollar interbank lending fell by the most since the fourth quarter of 2008.
  \item \textbf{Global resolution} (Gruenberg and Tucker). Of the world's 28 G-SIFIs, 12 are headquartered in the US or the UK. In a failure, the resolution authority would take control of the parent and wipe out shareholders. It would convert part of the unsecured debt into equity to recapitalise the group, and shrink, break up or close the businesses responsible.
\end{itemize}
\end{casestudy}

% ------------------------------------------------------------------
\section{Three world views}\label{sec:tv-views}

\paragraph{The money view: the present determines the present.} The course has focused on now: the day's clearing, with past investments and future promises fixed as
data. Cash outflows (discipline) meet cash inflows (elasticity); if outflows exceed inflows people borrow from the future and drive up the money rate of interest
\ts{22}{2}{0:45}. It is a flow view; stocks sit unchanged in the
background \ts{22}{2}{3:12}.

\paragraph{The economics view: the past determines the present.} Output today is a function of the capital stock accumulated by past investment, sunk and unchangeable
\ts{22}{2}{4:19}. Very much Adam Smith's view: worry about parsimony and profligacy, princes plundering the capital stock, slow
``turtle'' growth by accumulation \ts{22}{2}{5:21}. Through most of the history of economic thought the money view, natural
to bankers and central bankers, was a minority view against this one, which seems more natural and more ethical (the virtue of saving) \ts{22}{2}{6:24}.

\paragraph{The finance view: the future determines the present.} After World War~II a new view: the key variable is the \emph{price} of capital, the expected discounted
value of future cash flows, not the labour that built it. ``Radically subversive'': if capital is not productive it is worth zero; and an idea about the future, once
capitalized, gives purchasing power today, the tech bubble mentality \ts{22}{2}{7:49}.
Promises based on those dreams are promises to pay cash: back to the money view \ts{22}{2}{10:24}.

\begin{nfigure}
\begin{tikzpicture}[font=\scriptsize]
  \draw[thick,-{Stealth[length=4pt]}] (0,0) -- (13,0) node[right]{time};
  \fill (6.5,0) circle (2pt) node[below=3pt]{now};
  \node[align=center,draw,rounded corners,fill=defcol!10] at (2.2,1.1) {\textbf{economics view}\\past determines present\\quantity of capital (stock)};
  \node[align=center,draw,rounded corners,fill=keycol!10] at (6.5,1.1) {\textbf{money view}\\present determines present\\cash flows (flow)};
  \node[align=center,draw,rounded corners,fill=fincol!10] at (10.8,1.1) {\textbf{finance view}\\future determines present\\price of capital (stock)};
  \draw[-{Stealth}] (3.2,0.25) -- (6.2,0.25);
  \draw[-{Stealth}] (9.8,0.25) -- (6.8,0.25);
\end{tikzpicture}
\caption{Three views of the same elephant \ts{22}{2}{3:33}.}\label{fig:tv-views}
\end{nfigure}

Past investments generate cash inflows, the way they were financed generates cash outflows; dreams about the future require outlays greater than inflows; either someone
takes the other side or the dreams do not get their chance \ts{22}{2}{10:45}. Economics and finance are stock views and mirror images of each other;
the money view is a flow view \ts{22}{2}{11:25}. Mehrling's assertion: the big intellectual battle of the last 30 years has been economics versus finance,
settled by a modus vivendi (finance in the business school, economics in the economics department, each a minor subfield of the other); neither has much room for money.
His life project: make the money view talk to the finance view, as in the 19th century it talked to economics \ts{22}{2}{12:06}.

\begin{reading}[F.~Black, ``What a Non-Monetarist Thinks'']
The challenge that started Mehrling on this road: in a country with a smoothly working financial system, Black argued, the government cannot control the money stock;
money responds passively to the private sector's demand, and the Fed's monetary policy pronouncements are ``a charade'' with almost no effect on output, employment or
inflation. If true, ``we have just been wasting our time'' this semester \ts{22}{2}{13:27}.
\end{reading}

\begin{coursenote}
In class the piece is dated 1976 and said to be published. It was reprinted in F.~Black, \emph{Business Cycles and Equilibrium}, Basil Blackwell, 1987; the original
1976 date is as given in class and in Mehrling's biography of Black.
\end{coursenote}

% ------------------------------------------------------------------
\section{The economics view: commodity exchange}\label{sec:tv-econ}

The standard theory of value: apples and oranges, a production possibilities frontier, a budget line tangent to an indifference curve, its slope the relative price
$p_O/p_A$ \ts{22}{3}{0:01}. Two observations: there is no money (a barter relation between goods), and only a relative price is
determined, not the price level \ts{22}{3}{1:08}.

\begin{keyformulas}
\textbf{The quantity theory, pasted on} \ts{22}{3}{1:48}:
\[
  MV = p_O q_O + p_A q_A = PQ .
\]
Read left to right, money determines the price level; read right to left (the banking school, money as endogenous credit), transactions determine how much money is needed.
\end{keyformulas}

The velocity equation always had an uneasy relationship with the theory of value: an accounting identity turned into a theory, the closest the field came to science
\ts{22}{3}{3:11}. Irving Fisher's breakthrough: treat the two goods as consumption at dates 1 and 2, so the relative price involves $1/(1+r)$, and extend
the quantity equation to financial assets (the transactions version): more money means higher prices of goods and of assets, i.e.\ lower interest rates. From this tradition
came IS--LM: in the short run money affects output and interest rates, in the long run only the price level \ts{22}{3}{3:52}.

\begin{secondary}[Attempts to put money into general equilibrium]
The notes list three approaches, following Frank Hahn's complaint (\emph{Money and Inflation}, 1982) \mn{22}{3}\mn{22}{4}:
\begin{enumerate}
  \item overlapping generations, money as a store of value (N.~Wallace, ``The Overlapping Generations Model of Fiat Money'', 1980);
  \item search and matching (N.~Kiyotaki and R.~Wright, ``On Money as a Medium of Exchange'', 1989);
  \item Walrasian general equilibrium with transactions costs (R.~Starr, ``Monetary General Equilibrium with Transactions Costs'', 2005).
\end{enumerate}
All study commodity exchange with no financial asset but money; their results do not generalize to a world with financial assets, where money, paying no interest, is a dominated asset. The
most successful response to finance so far, Michael Woodford's \emph{Interest and Prices}, synthesizes the Taylor-rule consensus but looks only at goods prices, not asset prices: what remains is to do
for the modern model what Fisher did for the 19th-century one, link it to finance \mn{22}{8}.
\end{secondary}

\begin{intuition}[A patch]
The same structure extends to $n$ periods, intertemporal Arrow--Debreu general equilibrium. General equilibrium theorists, Frank Hahn above all, complained that it has no place
for money, and were unhappy with Fisher's patch. Post-war monetary economics can be told as attempts to squeeze money in (search theory and others); Mehrling bet that they
would not get far, and 25 years later is glad he chose finance instead \ts{22}{3}{6:58}.
\end{intuition}

% ------------------------------------------------------------------
\section{The finance view: risk}\label{sec:tv-finance}

\paragraph{CAPM.} (Say ``cap-M'': ``CAPM model'' says model twice.) Every capital asset is characterized by its expected return and its volatility; finance domesticates risk,
which general equilibrium, bringing in time without risk, did not; it is probabilistic risk, not true uncertainty, an assumption, but a step ahead
\ts{22}{4}{0:03}. Portfolios of securities form a frontier; investors like return and dislike
risk; the tangency of the line from the risk-free rate is the market portfolio, held by everyone, combined with more or less of the risk-free asset according to risk aversion;
beyond the tangency investors borrow \ts{22}{4}{2:10}.

\begin{keyformulas}
\textbf{CAPM} \ts{22}{4}{5:19}:
\begin{gather*}
  \text{security market line: } \E[R_i] = R_f + \beta_i\,(\E[R_m]-R_f),\\
  \text{capital market line: } \E[R_p] = R_f + (\E[R_m]-R_f)\,\frac{\sigma_p}{\sigma_m}.
\end{gather*}
\end{keyformulas}

\begin{definition}[New concepts: capital asset pricing model]
The \emph{CAPM} (Treynor, Sharpe, Lintner, Mossin, 1960s) prices each asset by its covariance with the market portfolio ($\beta_i$) times the market risk premium; ``capital
asset'' because it values the claims to capital \ts{22}{4}{0:03}.
\end{definition}

\paragraph{CAPM in balance sheets.} Two people with wealth 100 each: the risk-tolerant one holds 150 of the market portfolio and borrows 50; the risk-averse one holds 50 of
the market portfolio and 50 risk-free \ts{22}{4}{7:03}. Black's twist, not standard in teaching CAPM: the
50 lent and the 50 borrowed are the same 50, inside borrowing intermediated by a bank, loans on one side, deposits on the other \ts{22}{4}{8:48}.

\begin{nfigure}
\begin{taccts}
\tacct[1.9cm]{Risk tolerant}{Market portfolio \amt{150}}{Loan \amt{50}\\Net worth \amt{100}}\quad
\tacct[1.9cm]{Bank}{Loan \amt{50}}{Deposit \amt{50}}\quad
\tacct[1.9cm]{Risk averse}{Market portfolio \amt{50}\\Deposit \amt{50}}{Net worth \amt{100}}
\end{taccts}
\caption{Black's ``banking in a world without money'': the risk-free asset is a bank deposit funding a loan to the risk-tolerant investor \ts{22}{4}{7:45}.}\label{fig:tv-black}
\end{nfigure}

\begin{history}[Black's first paper]
Fischer Black's first publication was ``Banking and Interest Rates in a World without Money: The Effects of Uncontrolled Banking'', \emph{Journal of Bank Research} 1, 1970,
pp.~8--20, describing in words a world where banks intermediate between borrowers and lenders and money has no special role.\par
\emph{Reference:} P.~Mehrling, \emph{Fischer Black and the Revolutionary Idea of Finance}, Wiley, 2005.
\end{history}

\begin{reading}[Black, ``Banking and Interest Rates in a World Without Money'' (1970)]
The course file (\file{materials/readings/L22-World\_Without\_Money.pdf}) is the reprint as ch.~1 of Black's \emph{Business Cycles and Equilibrium} (1987), from the \emph{Journal of Bank Research}, Autumn 1970.
\begin{itemize}
  \item \textbf{The thought experiment.} Banks may offer checking accounts on any terms and pay interest on demand deposits; there are no reserve requirements. Then there is no reason for open market operations, and no reasonable definition of the quantity of money. Neither the quantity theory nor liquidity preference applies (p.~1).
  \item \textbf{Predecessors.} Black sides with Vickrey, Tobin, Gurley--Shaw and Patinkin: the price level is indeterminate. He rejects Johnson and Friedman--Schwartz, who predicted an unlimited rise in prices (pp.~1--4).
  \item \textbf{Accounts.} Every individual, business or government has one account that may be positive (a deposit) or negative (a loan) up to a maximum set by the bank. Deposits earn a ``standard wholesale money rate'', and an active inter-bank market evens out imbalances (pp.~4--6).
  \item \textbf{No government bonds.} The government simply borrows from the banks like everyone else. Its borrowing is limited by Congress; ``massive government spending that is not balanced by taxation would cause the financial system to break down'' (p.~6).
  \item \textbf{No quantity of money.} Any candidate (positive balances, net balances, unused credit lines) is either arbitrary or equal to bank capital or to total wealth. Hence ``the quantity of money cannot affect national income, employment, or the rate of inflation, because it does not exist'' (pp.~7--8).
  \item \textbf{Currency.} Currency is issued on demand and so is passive (p.~9).
  \item \textbf{Evolution of payment.} Black then builds the system up step by step: commodity money, then shares of a standard portfolio, then personal notes, then banks that guarantee notes (pp.~9--11).
\end{itemize}
\end{reading}

\paragraph{Two conclusions.} The bank intermediates between risk-averse and risk-tolerant; with diversification and limits to lending its loans might be nearly risk-free
\ts{22}{4}{9:34}:
\begin{enumerate}
  \item the risk-free rate is the rate that clears this loan market, and the risk premium is determined in the capital market too: no degree of freedom is left for the Fed
        to peg the short rate, hence ``a charade'' \ts{22}{4}{10:35};
  \item if the price of capital rises, both gain, the risk-tolerant three times as much: it wants to borrow more and buy more of the market portfolio, while the risk-averse
        wants to sell risky assets and lend more. The bank's balance sheet grows on both sides: the money supply must be endogenous, rising with asset prices, for asset
        markets to clear. Money is passive; read the quantity theory from right to left \ts{22}{4}{11:57}.
\end{enumerate}
This made Black the \emph{b\^ete noire} of monetarists and economists; he did not mind being a minority of one \ts{22}{4}{15:06}.

The notes complete Black's picture \mn{22}{8}\mn{22}{9}. Once people use banks for riskless borrowing and lending, they inevitably use them for payments too: a payment from a negative to a
positive account expands credit, from positive to negative contracts it, the elasticity of lecture~\ref{ch:clearinghouse}. But in the finance view that elasticity comes from the elasticity of
credit in capital markets, not the other way round: Black starts at the bottom of the hierarchy, the money view at the top, and both arrive at the intertwining of payments and credit. The big
difference: for Black the only constraint is wealth, allocated by risk tolerance in efficient markets; nothing about the survival constraint or dealers. (Black liked to say there is nothing we can
call the quantity of money, hence nothing to control \mn{22}{8}.)

% ------------------------------------------------------------------
\section{The education of Fischer Black}\label{sec:tv-education}

Black taught at Chicago and MIT, then joined Goldman Sachs, remarking that the world looks different from the banks of the Hudson than from the banks of the Charles
\ts{22}{5}{0:04}. As president of the American Finance Association he gave the address ``Noise'' at its New York meeting on 30 December 1985: an efficient market might be
defined as one in which price is ``within a factor of 2 of value'', and by that definition almost all markets are efficient almost all of the time (at least 90\%)
\ts{22}{5}{0:25}. A bombshell: prices need not be closely related to fundamentals \ts{22}{5}{1:07}.

\begin{intuition}[Treynor added to CAPM]
Black, who never took an economics or finance course, had been a strong believer in efficient markets, and Chicago hired him. Mehrling argues, in his biography, that the
``factor of two'' came from his friend Jack Treynor's dealer model: prices are set by dealers quoting bid--ask spreads, and when inventories are pushed far to one side prices can
deviate as far as the outside spread, without outsiders knowing unless they see dealers' inventories. Ironically Black learned CAPM itself from Treynor, whose correct
unpublished version preceded Sharpe's. The dealer model is the beginning of a bridge between finance and the money view \ts{22}{5}{1:28}.
\end{intuition}

\begin{history}[Black and Treynor]
Black (1938--1995), PhD in applied mathematics at Harvard, met Jack Treynor at Arthur D.~Little in the 1960s; Treynor's manuscript ``Toward a Theory of Market Value of Risky Assets''
(1962) anticipated the CAPM of W.~F.~Sharpe (1964). Black taught at Chicago (1971--75) and MIT (1975--84) and joined Goldman Sachs in 1984. ``Noise'' was published in the
\emph{Journal of Finance} 41(3), July 1986.\par
\emph{References:} F.~Black, ``Noise'', \emph{Journal of Finance} 41, 1986, 529--543; P.~Mehrling, \emph{Fischer Black and the Revolutionary Idea of Finance}, 2005.
\end{history}

In the notes, Black when he died still thought of the constraint on market liquidity as dealer \emph{capital}, not dealer \emph{funding}; the collapse of LTCM in 1998, whose relative-value strategy
he had warned about, ``loading up on liquidity risk'' and unable to refinance when liquidity dried up, taught the second lesson. ``The education of Fischer Black had proceeded through Lesson 2,
but not yet Lesson 3''; and Lesson 3 implies that ``monetary policy is not a charade. The answer to the challenge of finance is the money view'' \mn{22}{9}\mn{22}{10}.

By his death in 1995 Black had moved far from efficient markets, though not all the way to the money view; Mehrling imagines him in the audience: ``you came part of the way;
now the next logical step is here'' \ts{22}{5}{4:13}.

% ------------------------------------------------------------------
\section{Steps from the finance view to the money view}\label{sec:tv-steps}

Perfect liquidity, buying and selling at the efficient price, is an assumption that liquidity is a free good; in reality a dealer takes imbalances onto his balance sheet, and
trading itself moves prices, not just fundamentals \ts{22}{6}{0:03}.

\begin{keyformulas}
\textbf{Three steps} \ts{22}{6}{1:25}:
\begin{enumerate}
  \item market liquidity (buying or selling fast, in volume, without moving the price) depends on the dealer system: it is not free, and prices deviate from fundamentals;
  \item dealers' ability to make markets depends on their funding liquidity (and capital);
  \item the ultimate source of funding liquidity is the central bank.
\end{enumerate}
\end{keyformulas}

\paragraph{Funding, not just capital.} The finance literature studies how dealers' capital limits market making; but losing capital would not matter if dealers could borrow. The
money view stresses funding liquidity, which can vanish even if capital does not: ``I'm scared, I'm not going to fund you'', and the financing limits close in. Black surely got to
step one and probably to step two \ts{22}{6}{3:13}.

\paragraph{The third step,} the one Black would have resisted most, is a long way from ``money can only be passive''. But it is not all the way to IS--LM and fine-tuning: in a
financially developed economy markets and dealers play a large role in money market rates, and the central bank changes them essentially by being a dealer, taking pressure onto
its own balance sheet, which works because its liabilities are money, better than dealers' \ts{22}{6}{5:21}.
The hierarchy of money is antithetical to the finance view, where assets differ only by price and risk, ``all flat''; the money view is hierarchical, the best money, the next
best, then securities: where lectures~1 and~2 began \ts{22}{6}{7:06}.

% ------------------------------------------------------------------
\section{A money view of economics and finance}\label{sec:tv-moneyview}

Banks, or money dealers more generally, can relax the survival constraint of anyone who comes to them, but need not: they sell access to their balance sheet
\ts{22}{7}{0:04}. To the debtor bound by ``the dead hand of the past'' the banker may say ``pay me tomorrow, I will carry you'', taking a liquidity risk the debtor
cannot bear, or ``no way, you're out of business''; to the entrepreneur with a new app, needing cash outflows before any inflow, the banker may lend or refuse
\ts{22}{7}{0:45}.

\begin{keyformulas}
\textbf{Selective control over liquidity allocation} \ts{22}{7}{1:46}: the banking system decides who gets spending
power, who rolls over debt, who issues debt in the first place; by endorsing others' liabilities it turns them into means of payment. It chooses discipline or elasticity, for
individuals and the whole economy; the central bank influences it. This is where the money view connects with what economics cares about (investment) and what finance cares
about (asset prices).
\end{keyformulas}

Economics and finance abstract from money as the oil that greases transactions, and end up seriously incomplete; the money view does not contradict them, but ``the present
determines the present'' is a genuine aspect of the system, with real effects in crises and in normal times \ts{22}{7}{3:31}. It has always been a
minority view and may remain one, ``but if you want to understand money markets, it's the point of view you want to be grappling with'' \ts{22}{7}{4:34}.

\begin{intuition}[The future of banking]
The notes close with an image \mn{22}{10}\mn{22}{11}: the dead hand of the past and the ephemeral dream of the future meet in the survival constraint of the present. Banks let agents with obligations
they cannot now meet put off the problem, and agents with brilliant plans they cannot now finance issue obligations; in doing so they make payments the agents cannot (yet) make. ``The essential
core of the banking function is the selective bearing of liquidity risk'': as long as there is liquidity risk, banks have a crucial allocative role. Regulation is about the externalities, the
ways the private profit motive in bearing liquidity risk deviates from the social good: the more risk banks bear, the more they charge for more, until they will bear no more and the system breaks
down. Central banks can help as lender and dealer of last resort, and so also by intervening earlier, when private and social begin to diverge.
\end{intuition}

% ------------------------------------------------------------------
\flushside
\section*{Key formulas and ideas}
\begin{keyformulas}
\begin{tabularx}{\linewidth}{@{}l L@{}}
Three views & economics: past $\to$ present (quantity of capital); finance: future $\to$ present (price of capital); money: present $\to$ present (cash flows)\\
Economics & relative prices from GE; $MV=PQ$ pasted on; Fisher $\to$ IS--LM\\
CAPM & $\E[R_i]=R_f+\beta_i(\E[R_m]-R_f)$; Black: risk-free asset $=$ bank deposit, money endogenous and passive\\
Noise & price within a factor of 2 of value: the dealer's outside spread\\
Three steps & market liquidity $\leftarrow$ dealers $\leftarrow$ funding liquidity $\leftarrow$ central bank\\
Money view & hierarchy of money; banks' selective control over liquidity allocation\\
\end{tabularx}
\end{keyformulas}

\section*{Exercises}

\begin{exercise}[Black's banking world]
In the two-person example the price of the market portfolio rises by 20\%. Compute both net worths, and the new loans and deposits if each restores its original ratio of
risky assets to net worth. By how much does the money supply grow?
\end{exercise}

\begin{exercise}[Reading the quantity equation]
Explain the difference between reading $MV=PQ$ from left to right and from right to left, and which reading fits Black's world, the banking school and the money view.
\end{exercise}

\begin{exercise}[A factor of two]
Using Treynor's model, explain how a dealer's price can deviate from value as far as the outside spread, and what determines the width of that spread in an equity market and
in a money market.
\end{exercise}

\begin{exercise}[Three views of one event]
Describe the collapse of AIG in September 2008 from the economics view, the finance view and the money view. Which features does each view see, and which does it miss?
\end{exercise}
